% !TEX root = ../Main.tex
\chapter{平面向量基础知识复习}

\section{坐标与共线}

点 $A(x,y)$ 对应位置向量 $\vec{OA}=(x,y)$。取一组基 $\vec e_1,\vec e_2$（且 $\vec e_1\perp\vec e_2$），把 $\vec{OA}$ 沿基展开
\[
\vec{OA}=x\vec e_1+y\vec e_2\ \Longrightarrow\ \vec{OA}=(x,y).
\]

\textbf{共线}：$\vec a,\vec b$ 共线，就是存在 $\lambda$ 使
\[
\vec a=\lambda\vec b\quad(\vec b\ne\vec 0).
\]

\state{共线的推论：三点共线}{
    设 $A,B,C$ 三点共线，在外面任取一点 $O$ 当参考点，则
    \[
    \vec{OA}=\lambda\,\vec{OB}+\mu\,\vec{OC},\qquad \lambda+\mu=1.
    \]
}

\section{平面向量基本定理}

\state{平面向量基本定理}{
    $\vec e_1$ 与 $\vec e_2$ 不共线，则在 $\vec e_1,\vec e_2$ 所在平面内的任意一个向量 $\vec m$，都可以表示为
    \[
    \vec m=x\vec e_1+y\vec e_2.
    \]
    特别地，当 $|\vec e_1|=|\vec e_2|=1$ 且 $\vec e_1\cdot\vec e_2=0$ 时，就建立起平面直角坐标系 $xOy$，此时
    \[
    \vec m=(x,y).
    \]
}

“不共线”是关键——共线的两个向量只能张成一条直线，张不出整个平面。

\section{点乘}

\state{点乘的定义与本质}{
    \[
    \vec a\cdot\vec b=|\vec a|\,|\vec b|\cos\langle\vec a,\vec b\rangle.
    \]
    本质是\textbf{“投影 $\times$ 模长”}：$|\vec a|\cos\langle\vec a,\vec b\rangle$ 是 $\vec a$ 在 $\vec b$ 方向上的投影，再乘上 $|\vec b|$。
}

\begin{center}
\begin{tikzpicture}[>=Stealth, scale=1]
    % 锐角情形：投影 b' 与 a 同向
    \draw[->, thick] (0,0) -- (3,0) node[right] {$\vec a$};
    \draw[->, thick] (0,0) -- (1.4,1.1) node[above] {$\vec b$};
    \draw[dashed] (1.4,1.1) -- (1.4,0);
    \draw[->, thick, blue] (0,0) -- (1.4,0);
    \node[below, blue] at (0.7,-0.05) {\footnotesize $\vec b'$};
\end{tikzpicture}
\qquad\qquad
\begin{tikzpicture}[>=Stealth, scale=1]
    % 钝角情形：投影 b' 与 a 反向
    \draw[->, thick] (0,0) -- (3,0) node[right] {$\vec a$};
    \draw[->, thick] (0,0) -- (-1.1,1.1) node[above left] {$\vec b$};
    \draw[dashed] (-1.1,1.1) -- (-1.1,0);
    \draw[->, thick, blue] (0,0) -- (-1.1,0);
    \node[below, blue] at (-0.55,-0.05) {\footnotesize $\vec b'$};
\end{tikzpicture}
\end{center}

\section{坐标形式：点乘、模长、夹角}

设 $\vec{OA}=x_1\vec e_1+y_1\vec e_2$，$\vec{OB}=x_2\vec e_1+y_2\vec e_2$，其中 $\vec e_1,\vec e_2$ 是单位正交基（$\vec e_1^2=\vec e_2^2=1$，$\vec e_1\cdot\vec e_2=0$）。直接展开：
\[
\vec{OA}\cdot\vec{OB}=x_1x_2\,\vec e_1^2+(x_1y_2+x_2y_1)\,\vec e_1\cdot\vec e_2+y_1y_2\,\vec e_2^2=x_1x_2+y_1y_2.
\]

于是对 $\vec a=(x_1,y_1),\ \vec b=(x_2,y_2)$：
\[
\vec a\cdot\vec b=x_1x_2+y_1y_2,\qquad
|\vec a|=\sqrt{\vec a^2}=\sqrt{x_1^2+y_1^2},\quad |\vec b|=\sqrt{x_2^2+y_2^2},
\]
\[
\cos\langle\vec a,\vec b\rangle=\frac{\vec a\cdot\vec b}{|\vec a|\,|\vec b|}=\frac{x_1x_2+y_1y_2}{\sqrt{x_1^2+y_1^2}\,\sqrt{x_2^2+y_2^2}}.
\]

\section{平行与垂直}

\state{平行 $\equiv$ 共线}{
    $\vec a=(x_1,y_1),\ \vec b=(x_2,y_2)$ 平行（共线），从比例上看是
    \[
    \frac{x_1}{x_2}=\frac{y_1}{y_2}\quad\text{或}\quad \frac{x_1}{y_1}=\frac{x_2}{y_2},
    \]
    但分母可能为零、不够严谨。\textbf{标准写法}是：存在 $\lambda$ 使
    \[
    \begin{cases}x_1=\lambda x_2,\\ y_1=\lambda y_2.\end{cases}
    \]
}

\textbf{垂直}：$\vec a\perp\vec b\iff\vec a\cdot\vec b=x_1x_2+y_1y_2=0$。

\section{升维：从平面到空间}

上面这套结论，到了空间里\textbf{只是多了一个坐标}：向量写成 $\vec a=(x,y,z)$，点乘多一项 $z_1z_2$，模长多一项 $z^2$，平行、垂直、夹角的形式完全不变。
